% !TeX TS-program = lualatex

\documentclass[aspectratio=169,13pt,t] {beamer}
%
% Choose how your presentation looks.
%
% For more themes, color themes and font themes, see:
% http://deic.uab.es/~iblanes/beamer_gallery/index_by_theme.html
%
\mode<presentation>
{
  \usetheme{Madrid}      % or try Darmstadt, Madrid, Warsaw, ...
  \usecolortheme{beaver} % or try albatross, beaver, crane, ...
  \usefonttheme{serif}  % or try serif, structurebold, ...
%  \setbeamertemplate{navigation symbols}{}
  \setbeamertemplate{caption}[numbered]
  \setbeamertemplate{items}[default]
} 

\usepackage[,UTF8,linespread=1.2]{ctex}
\usepackage{wrapfig}
\usepackage{caption}
\captionsetup[figure]{name=图,labelsep=period,}

\usepackage[,most]{tcolorbox}

\usepackage{picinpar}
\usepackage{pgfplots}
\pgfplotsset{
	compat=1.18,
	every axis legend/.append style={
		at={(0.8,0.6)},
		anchor=north west,
		legend cell align=left,
		draw = red!75!black,
	},
}

\usepackage{tkz-euclide}
\usepackage{tkz-elements}
\usepackage{xcolor}
\definecolor{xy_blue}{cmyk}{0.6,0.35,0.0,0.4}
\definecolor{xy_orange}{cmyk}{0.0,0.7,1.0,0.0}
\definecolor{xy_green}{cmyk}{0.2,0.0,1.0,0.15}
\xdefinecolor{xy_gold}{RGB}{135,103,79}
\xdefinecolor{xy2_gold}{cmyk}{0,0.35,0.75,0.05}


\setbeamercolor{section number projected}{fg=red,bg=red!20}%控制球的颜色
\setbeamercolor{subsection number projected}{fg=blue,bg=blue!40}
\setbeamercolor{item}{fg=xy_gold} %设置列表项颜色

\usepackage{amsthm}
\newtheorem{peroration}{结论：}
\newtheorem{exploration}{探究：}
\newtheorem{prooft}{证明：}
\newtheorem{question}{问题}
\newtheorem*{solve}{解:} 
\newtheorem{instance}{例}   % 解：独立计数器
\newtheorem{exercise}{练习}
\BeforeBeginEnvironment{peroration}{%
	\setbeamercolor{block title}{bg=xy_green,fg=white}
	\setbeamercolor{block body}{bg=xy_green!10}
}
\AfterEndEnvironment{peroration}{
	\setbeamercolor{block title}{bg=xy_blue,fg=white}
	\setbeamercolor{block body}{bg=xy_blue!10}
}
\BeforeBeginEnvironment{exploration}{%
	\setbeamercolor{block title}{bg=xy_orange,fg=white}
	\setbeamercolor{block body}{bg=xy_orange!10}
}
\AfterEndEnvironment{peroration}{
	\setbeamercolor{block title}{bg=xy_blue,fg=white}
	\setbeamercolor{block body}{bg=xy_blue!10}
}
\BeforeBeginEnvironment{solution}{%
	\setbeamercolor{block title}{bg=xy_green,fg=white}
	\setbeamercolor{block body}{bg=xy_green!10}
}
\AfterEndEnvironment{solution}{
	\setbeamercolor{block title}{bg=xy_green,fg=white}
	\setbeamercolor{block body}{bg=xy_green!10}
}
\BeforeBeginEnvironment{prooft}{%
	\setbeamercolor{block title}{bg=xy_green,fg=white}
	\setbeamercolor{block body}{bg=xy_green!10}
}
\AfterEndEnvironment{prooft}{
	\setbeamercolor{block title}{bg=xy_green,fg=white}
	\setbeamercolor{block body}{bg=xy_green!10}
}
\BeforeBeginEnvironment{question}{%
	\setbeamercolor{block title}{bg=xy_orange,fg=white}
	\setbeamercolor{block body}{bg=xy_orange!10}
}
\AfterEndEnvironment{question}{
	\setbeamercolor{block title}{bg=xy_green,fg=white}
	\setbeamercolor{block body}{bg=xy_green!10}
}
\BeforeBeginEnvironment{exercise}{%练习
	\setbeamercolor{block title}{bg=xy_green,fg=white}
	\setbeamercolor{block body}{bg=xy_green!10}
}
\AfterEndEnvironment{exercise}{%练习
	\setbeamercolor{block title}{bg=xy_green,fg=white}
	\setbeamercolor{block body}{bg=xy_green!10}
}
\BeforeBeginEnvironment{solve}{%解：
	\begingroup
	\setbeamercolor{block title}{bg=xy_orange!60,fg=white}
	\setbeamercolor{block body}{bg=xy_orange!10}
}
\AfterEndEnvironment{solve}{%解：
	\endgroup
}

\setbeamercolor{prooft}{fg=red, bg=green!10}       % 标题颜色
\setbeamercolor{prooft body}{fg=red, bg=red!5}  % 内容块背景
\setbeamerfont{prooft head}{series=\bfseries}       % 标题字体
\setbeamerfont{question body}{shape=\upshape}         % 内容字体
\renewcommand{\thefigure}{\arabic{figure}}

\newtcolorbox{summarybox}[1][]{
	colback=purple!10,
	colframe=purple!70!black,
	fonttitle=\bfseries,
	title=#1,
	arc=1mm,
	boxrule=0pt,drop shadow,
	left=2mm, right=2mm, top=2mm, bottom=2mm,
}
\newtcolorbox{summarybox-answer}{
	width=0.9\textwidth,
	colback=white,
	blanker,
	borderline west={2pt}{20pt}{cyan},
	left=10mm, right=2mm, top=2mm, bottom=2mm,
	}
\title[3.1.2函数的表示法]{3.1.2函数的表示法}
\subtitle{解析法、列表法、图象法}
\author[好好学习，努力向上]{}
\institute{}
\date{2026.10.12}
%\setbeamertemplate{theorems}[numbered]
\theoremstyle{definition}
\newtheorem{exam}{例}
\AtBeginSection[]
{
  \begin{frame}<beamer>
    \frametitle{本节主要内容}
    \tableofcontents[currentsection,currentsubsection]
  \end{frame}
}

\AtBeginDocument{
	\setlength{\abovedisplayskip}{3pt}
	\setlength{\belowdisplayskip}{3pt}
	\setbeamertemplate{theorems}[numbered]
}

\usepackage{pgfplots}
\pgfplotsset{compat=1.18}
\usetikzlibrary{arrows.meta}
\tikzset{>=Stealth}

\begin{document}
\begin{frame}{}
	\titlepage
\end{frame}

\begin{frame}{内容概览}
	\tableofcontents
\end{frame}

\section{函数的三种表示法}

\begin{frame}{函数的三种表示法}
	我们在初中已经接触过函数的三种表示法：
	\begin{columns}
		\column{0.32\textwidth}
		\begin{summarybox}[解析法]
			用解析式表示两个变量之间的对应关系.
		\end{summarybox}
		\column{0.32\textwidth}
		\begin{summarybox}[列表法]
			列出表格来表示两个变量之间的对应关系.
		\end{summarybox}
		\column{0.32\textwidth}
		\begin{summarybox}[图象法]
			用图象表示两个变量之间的对应关系.
		\end{summarybox}
	\end{columns}
	\vspace{1em}
	\begin{summarybox-answer}
		这三种方法是常用的函数表示法.
	\end{summarybox-answer}
\end{frame}

\section{例4：三种表示法}

\begin{frame}%{例4：笔记本问题}
	\setcounter{exam}{3}
	
	\begin{exam}
	某种笔记本的单价是 5 元，买 \(x\)（\(x \in \{1, 2, 3, 4, 5\}\)）个笔记本需要 \(y\) 元. 试用函数的三种表示法表示函数 \(y = f(x)\).
	\end{exam}
	\pause
	\begin{solve}
		这个函数的定义域是数集 \(\{1, 2, 3, 4, 5\}\).
	\end{solve}
\end{frame}

\begin{frame}%{例4（1）解析法}
	\begin{solve}
		用解析法可将函数 \(y = f(x)\) 表示为
		\[
			y = 5x,\quad x \in \{1, 2, 3, 4, 5\}.
		\]
	\end{solve}
	\vspace{-1em}
	\begin{summarybox}%[列表法]
		用列表法可将函数 \(y = f(x)\) 表示为：
	\end{summarybox}
	\begin{center}
		\renewcommand{\arraystretch}{1.4}
		\begin{tabular}{|c|c|c|c|c|c|}
			\hline
			笔记本数 \(x\) & 1 & 2 & 3 & 4 & 5 \\ \hline
			钱数 \(y\) & 5 & 10 & 15 & 20 & 25 \\ \hline
		\end{tabular}
	\end{center}
\end{frame}

\begin{frame}%{例4（2）图象法}
	用图象法可将函数 \(y = f(x)\) 表示为图 \ref{fig:img1}.
	
	\vspace{1em}
	\begin{columns}[t]
		\column{0.5\textwidth}
		\begin{tikzpicture}[baseline={(0,5)}]
			\begin{axis}[
				width=7cm, height=6cm,
				xmin=0, xmax=5.5,
				ymin=0, ymax=27,
				xtick={1,2,3,4,5},
				ytick={5,10,15,20,25},
				axis lines=middle,
				axis line style={->},
				xlabel={$x$}, ylabel={$y$},
				xlabel style={below right},
				ylabel style={above left},
				tick label style={font=\small},
			]
			\addplot[only marks, mark=*, mark size=3pt, color=magenta!70!purple] coordinates {
				(1,5) (2,10) (3,15) (4,20) (5,25)
			};
			\end{axis}
		\end{tikzpicture}
		\captionof{figure}{}
		\label{fig:img1}
		
		\column{0.45\textwidth}
%		\vspace{-2em}
		\begin{summarybox-answer}
			函数图象既可以是连续的曲线，也可以是直线、折线、离散的点等.
			
			~
			
			\textbf{思考：}判断一个图形是不是函数图象的依据是什么？
		\end{summarybox-answer}
	\end{columns}
\end{frame}

\begin{frame}%{思考}
	\begin{summarybox}[思考]
		\begin{enumerate}[(1)]
			\item 比较函数的三种表示法，它们各自的特点是什么？
			\item 所有函数都能用解析法表示吗？请你举出实例加以说明.
		\end{enumerate}
	\end{summarybox}
	\vspace{1em}
	\pause
	\begin{summarybox-answer}
		\textbf{解析法}：精确、简洁，便于研究性质；\\
		\textbf{列表法}：直观，直接看出对应值；\\
		\textbf{图象法}：形象，直观看出变化趋势.
	\end{summarybox-answer}
\end{frame}

\section{例5：分段函数}

\begin{frame}{例5：画出函数 \(y = |x|\) 的图象}
%	\begin{exam}
%	画出函数 \(y = |x|\) 的图象.
%	\end{exam}
%	\pause
	\begin{solve}
		由绝对值的概念，我们有
		\[
			y = \begin{cases} -x, & x < 0, \\ x, & x \geqslant 0. \end{cases}
		\]
		所以，函数 \(y = |x|\) 的图象如图 \ref{fig:img2} 所示.
	\end{solve}
\end{frame}

\begin{frame}{图 \ref{fig:img2}：\(y = |x|\) 的图象}
	\begin{columns}[t]
		\column{0.55\textwidth}
		\begin{tikzpicture}[baseline=5cm]
			\begin{axis}[
				width=8cm, height=6cm,
				xmin=-3.5, xmax=3.5,
				ymin=-0.5, ymax=5.5,
				xtick={-3,-2,-1,0,1,2,3},
				ytick={1,2,3,4,5},
				axis lines=middle,
				axis line style={->},
				xlabel={$x$}, ylabel={$y$},
				xlabel style={below right},
				ylabel style={above left},
				tick label style={font=\small},
				tick align=inside,% 刻度线位置：xtick align= ,ytick align= ,inside,outside,center,默认为center.
				tick style={line width=.4pt,},	%刻度线样式，线宽，颜色等
				major tick length=2pt, %主刻度线长度 minor tick length={⟨dimension⟩}次刻度线长度
			]
			\addplot[color=pink!80!purple, line width=1.5pt, domain=-3:0] {abs(x)};
			\addplot[color=pink!80!purple, line width=1.5pt, domain=0:3] {x};
			\end{axis}
		\end{tikzpicture}
		\captionof{figure}{}
		\label{fig:img2}
		
		\column{0.4\textwidth}
%		\vspace{3em}
		\begin{summarybox}%[分段函数]
			像例 5 中 \(y = \begin{cases} -x, & x < 0, \\ x, & x \geqslant 0 \end{cases}\) 这样的函数称为\textbf{分段函数}.
		\end{summarybox}
%		\vspace{1em}
		\begin{summarybox-answer}
			生活中，有很多可以用分段函数描述的实际问题，如出租车的计费、个人所得税纳税额等.
		\end{summarybox-answer}
	\end{columns}
\end{frame}

\section{例6：max函数}

\begin{frame}%{例6：求 \(M(x) = \max\{f(x), g(x)\}\)}
	\setcounter{exam}{5}
	\begin{exam}
	给定函数 \(f(x) = x + 1\)，\(g(x) = (x + 1)^2\)，\(x \in \mathbf{R}\)，
	\begin{enumerate}[(1)]
		\item 在同一直角坐标系中画出函数 \(f(x)\)，\(g(x)\) 的图象；
		\item \(\forall x \in \mathbf{R}\)，用 \(M(x)\) 表示 \(f(x)\)，\(g(x)\) 中的最大者，记为
		\[
			M(x) = \max\{f(x),\, g(x)\}.
		\]
		例如，当 \(x = 2\) 时，\(M(2) = \max\{f(2),\, g(2)\} = \max\{3,\, 9\} = 9\).
		
		请分别用图象法和解析法表示函数 \(M(x)\).
	\end{enumerate}
	\end{exam}
\end{frame}

\begin{frame}%{例6（1）：\(f(x)\) 和 \(g(x)\) 的图象}
	\vspace{-1em}
	\begin{solve}
	(1)	在同一直角坐标系中画出函数 \(f(x)\)，\(g(x)\) 的图象（图 \ref{fig:img3}）.
	\end{solve}
\begin{center}
		\begin{tikzpicture}
			\begin{axis}[
				width=8cm, height=6cm,
				xmin=-5.5, xmax=5.5,
				ymin=-2.5, ymax=5.5,
				xtick={-5,-4,-3,-2,-1,0,1,2,3,4,5},
				ytick={-2,-1,0,1,2,3,4,5},
				axis lines=middle,
				axis line style={->},
				xlabel={$x$}, ylabel={$y$},
				xlabel style={below right},
				ylabel style={above left},
				tick label style={font=\tiny},
			]
			% g(x) = (x+1)^2
			\addplot[color=pink!80!purple, line width=1.5pt, domain=-3.2:2] {(x+1)^2};
			\node[above , pink!30!black] at (-3, 4.5) {$g(x)=(x+1)^2$};
			% f(x) = x+1
			\addplot[color=cyan!70!blue, line width=1.5pt, domain=-3.5:3] {x+1};
			\node[below right, cyan!70!blue] at (2.5, 4) {$f(x)\!=\!x\!+\!1$};
			\end{axis}
		\end{tikzpicture}
		\captionof{figure}{}
		\label{fig:img3}
		\end{center}
		
	
\end{frame}

\begin{frame}%{例6（2）：\(M(x)\) 的图象}
	\vspace{-1em}
	\begin{solve}
		由图 \ref{fig:img3} 中函数取值的情况，结合函数 \(M(x)\) 的定义，可得函数 \(M(x)\) 的图象（图 \ref{fig:img4}）.
		\end{solve}
\begin{center}
		\begin{tikzpicture}
			\begin{axis}[
				width=7.5cm, height=6cm,
				xmin=-5.5, xmax=5.5,
				ymin=-2.5, ymax=5.5,
				xtick={-5,-4,-3,-2,-1,0,1,2,3,4,5},
				ytick={-2,-1,0,1,2,3,4,5},
				axis lines=middle,
				axis line style={->},
				xlabel={$x$}, ylabel={$y$},
				xlabel style={below right},
				ylabel style={above left},
				tick label style={font=\tiny},
			]
			% M(x) = max of (x+1)^2 and x+1
			% 交点：(x+1)^2 = x+1 → (x+1)x = 0 → x = -1 or x = 0
			% x <= -1: g(x) = (x+1)^2 (抛物线在上)
			% -1 < x <= 0: f(x) = x+1 (直线在上)
			% x > 0: g(x) = (x+1)^2 (抛物线在上)
			\addplot[color=magenta!80!purple, line width=1pt, domain=-4:-1] {(x+1)^2};
			\addplot[color=green!50!black, line width=1pt, domain=-1:0] {x+1};
			\addplot[color=magenta!80!purple, line width=1pt, domain=0:2] {(x+1)^2};
			\node[above right, magenta!80!purple] at (1.5, 4.5) {$M(x)$};
			\end{axis}
		\end{tikzpicture}
		\captionof{figure}{}
		\label{fig:img4}
		\end{center}
		

\end{frame}

\begin{frame}%{例6（2）：\(M(x)\) 的解析式}
	\begin{solve}
		由 \((x+1)^2 = x + 1\)，得 \(x(x+1) = 0\).
		
		解得 \(x = -1\)，或 \(x = 0\).
	\end{solve}
	\vspace{1em}
	\pause
	结合图 \ref{fig:img3}，得出函数 \(M(x)\) 的解析式为：
	\[
		M(x) = \begin{cases}
			(x+1)^2, & x \leqslant -1, \\
			x+1, & -1 < x \leqslant 0, \\
			(x+1)^2, & x > 0.
		\end{cases}
	\]
	\vspace{1em}
	\pause
	\begin{summarybox-answer}
		\textbf{思考：}你能用其他方法求出 \(M(x)\) 的解析式吗？
	\end{summarybox-answer}
\end{frame}

\begin{frame}{课堂小结}
	
	本节课学习的内容：
	
	\begin{enumerate}[(1)]
		\item 函数的3种表示法
		\item 分段函数
	\end{enumerate}
	
\end{frame}

\begin{frame}{布置作业}
	\begin{summarybox}[课后作业]
		\begin{enumerate}
			\item 教材 P69 练习 1、2、3
			\item P72 习题3.1\quad 3、4
		\end{enumerate}
	\end{summarybox}
\end{frame}

\section{课堂练习}

\begin{frame}{练习 1：圆内接矩形}
	\begin{exercise}
	如图，把直截面半径为 25\,cm 的圆柱形木头锯成直截面为矩形的木料，如果矩形的一边长为 \(x\)（单位：cm），面积为 \(y\)（单位：\(\text{cm}^2\)），把 \(y\) 表示为 \(x\) 的函数.
	\end{exercise}
	\pause
	\vspace{-1em}
	\begin{columns}[t]
		\column{0.65\textwidth}
	
	\begin{solve}
		设矩形另一边长为 \(z\)，由勾股定理：
		\[
			x^2 + z^2 = 50^2 = 2500 \quad\Rightarrow\quad z = \sqrt{2500 - x^2}.
		\]
		所以面积为：
		\[
			y = x \cdot z = x\sqrt{2500 - x^2},\quad x \in (0, 50).
		\]
	\end{solve}
	
	\column{0.35\textwidth}
	\begin{center}
		
	
	\directlua{	
	init_elements()
	z.A = point(0,0)
	z.B = point(3,0)
	z.C = point(3,2)
	z.D = point(0,2)
	z.O = tkz.midpoint(z.A,z.C)	
	}
\begin{tikzpicture}[baseline=1cm]
	\tkzGetNodes		
	\tkzDrawCircle[thick,black,fill=orange!20](O,A)
	\tkzDrawPolygons[thick,fill=orange!40!black!60](A,B,C,D)
	\tkzDrawPoints[shape=circle,fill=black,size=0.2](O)	
	\tkzDrawSegment[](O,C)
	\tkzLabelSegment[sloped,above left=0pt](O,C){25cm}
	\tkzLabelSegment(A,B){$x$}
\end{tikzpicture}
\end{center}
\end{columns}
\end{frame}

\begin{frame}{练习 2：画出 \(y = |x - 2|\) 的图象}
	\begin{exercise}
	画出函数 \(y = |x - 2|\) 的图象.
	\end{exercise}
	\pause
	\begin{columns}[t]
		\column{0.45\textwidth}
	\begin{solve}
		由绝对值的概念：
		\[
			y = \begin{cases} -(x-2) = 2-x, & x < 2, \\ x-2, & x \geqslant 2. \end{cases}
		\]
		图象是顶点在 \((2, 0)\) 的 V 形，开口向上.
	\end{solve}
	\column{0.45\textheight}
	\begin{tikzpicture}[baseline=3.5cm]
		\begin{axis}[
			width=5cm, height=5cm,
			axis lines=middle,
			xlabel=$x$, ylabel=$y$,
			xmin=-2.2, xmax=6.2, ymin=-1, ymax=5.2,
			xtick={2},
			ytick={2},
			grid=major, grid style={gray!20},
			samples=200,
			]
			\addplot[thick,blue,domain=-2:6] {abs(x-2)};
			\addplot[mark=*,mark size=2pt,only marks] coordinates {(2,0)};
%			\node[below right] at (axis cs:2,0) {$(2,0)$};
		\end{axis}
	\end{tikzpicture}
	
\end{columns}
\end{frame}

\begin{frame}{练习 3：\(m(x) = \min\{f(x), g(x)\}\)}
	\vspace{-1em}
	\begin{exercise}
	给定函数 \(f(x) = -x + 1\)，\(g(x) = (x-1)^2\)，\(x \in \mathbf{R}\)，
	\begin{enumerate}[(1)]
		\item 画出函数 \(f(x)\)，\(g(x)\) 的图象；
		\item \(\forall x \in \mathbf{R}\)，用 \(m(x)\) 表示 \(f(x)\)，\(g(x)\) 中的最小者，记为 \(m(x) = \min\{f(x), g(x)\}\)，请分别用图象法和解析法表示函数 \(m(x)\).
	\end{enumerate}
	\end{exercise}
	\pause
	\begin{summarybox-answer}
		\textbf{提示：}先求交点 \(-x+1 = (x-1)^2\)，再分段讨论.
		
		\textbf{答案：}$m(x)=\begin{cases}
			-x+1,& x \leqslant 0,\\
			(x-1)^2,& 0 <x \leqslant 1,\\
			-x+1,& x>1
		\end{cases}$
	\end{summarybox-answer}
\end{frame}

\begin{frame}{图像}
	\begin{columns}
		
	\column{0.45\textwidth}
	\begin{tikzpicture}
		\begin{axis}[
			width=6.5cm, height=6.5cm,
			axis lines=middle,
			xlabel=$x$, ylabel=$y$,
			xmin=-2, xmax=4, ymin=-1, ymax=4,
			xtick={-1,0,1,2,3},
			ytick={1,2,3},
			grid=major, grid style={gray!20},
%			legend pos=outer north east,
			samples=100,
			]
			\addplot[very thick,blue,domain=-2:4] {-x+1};
			\addlegendentry{$f(x)=-x+1$}
			\addplot[very thick,red,domain=-2:4] {(x-1)^2};
			\addlegendentry{$g(x)=(x-1)^2$}
			\addplot[only marks,mark=*,mark size=2pt] coordinates {(0,1)(1,0)};
			
		\end{axis}
	\end{tikzpicture}
	
	\column{0.45\textwidth}
	\begin{tikzpicture}
		\begin{axis}[
			width=6.5cm, height=6.5cm,
			axis lines=middle,
			xlabel=$x$, ylabel=$y$,
			xmin=-2, xmax=4, ymin=-1, ymax=4,
			xtick={-1,0,1,2,3},
			ytick={1,2,3},
			grid=major, grid style={gray!20},
			samples=100,
			]
			% x≤0段：m(x)=f(x)=-x+1
			\addplot[very thick,blue,domain=-2:0] {-x+1};
			% 0≤x≤1段：m(x)=g(x)=(x-1)²
			\addplot[very thick,red,domain=0:1] {(x-1)^2};
			% x≥1段：m(x)=f(x)=-x+1
			\addplot[very thick,blue,domain=1:4] {-x+1};
			% 分段点
			\addplot[only marks,mark=*,mark size=2pt] coordinates {(0,1)(1,0)};
			\node at (1,1) {$m(x)$};
		\end{axis}
	\end{tikzpicture}	
\end{columns}
\end{frame}

\end{document}
